Fix an integer $k\ge 1$. For each prime $p$, define $S_{k,p}\subset\mathbb{N}$ by
\[
S_{k,p}:=\{n\in\mathbb{N}:\ \text{if }p_1<p_2<\cdots<p_{\omega(n)}\text{ are the distinct prime divisors of }n,\text{ then }\omega(n)\ge k\text{ and }p_k=p\}.
\]
Define
\[
d_k(p):=\lim_{x\to\infty}\frac{1}{x}\#\{n\le x: n\in S_{k,p}\},
\]
provided the limit exists (natural density). The question asks whether, for fixed $k$, the function $p\mapsto d_k(p)$ is \emph{unimodular/unimodal} along the primes: i.e. whether there exists a prime $p^*(k)$ such that $d_k(p)$ is nondecreasing for primes $p\le p^*(k)$ and nonincreasing for primes $p\ge p^*(k)$. 

\medskip
\noindent\textbf{Ambiguity note.} The problem is phrased as a question. A minimal “yes/no” statement consistent with standard conventions is:

\smallskip
\noindent\emph{(Universal unimodality claim)}: “For every fixed $k\ge 1$, the map $p\mapsto d_k(p)$ is unimodular.”

\noindent Below we give an explicit counterexample (already for $k=4$), so this universal claim is false.
